\section{Bicentralizers and the reduction}\label{sec:bicentralizers}

We recall the definition to fix both the topology and the quantifiers.
For a faithful normal state $\phi$, put
$\norm{x}_\phi=\phi(x^*x)^{1/2}$.
The bicentralizer consists of those $a\in M$ for which, for every
$\varepsilon>0$, there is $\delta>0$ such that
\[
  u\in\mathcal U(M),\quad
  \norm{u\phi-\phi u}_{M_*}<\delta
  \quad\Longrightarrow\quad
  \norm{u^*au-a}_\phi<\varepsilon.
\]
This is Haagerup's unitary formulation of the usual bicentralizer
\cite[proof of Lemma~1.2]{Haagerup1987}.
Equivalently, its elements asymptotically commute strongly with every
uniformly bounded sequence that asymptotically centralizes $\phi$ in
predual norm. We use established structural results in this equivalent
form.

\begin{proposition}[Established bicentralizer results]\label{bc:inputs}
Let $M$ be a type $\mathrm{III}_1$ factor with separable predual and
$\phi$ a faithful normal state.
\begin{enumerate}
\item\label{bc:reduction}
If $B=\BC(M,\phi)$ is nonscalar, then $B$ is a type
$\mathrm{III}_1$ subfactor with a normal $\phi$-preserving conditional
expectation, and
\[
  \BC(B,\phi|_B)=B.
\]
\item\label{bc:flow}
There is a continuous $\phi$-preserving action
$\beta^\phi:\R_+^*\curvearrowright\BC(M,\phi)$ characterized by
\[
  \norm{w_n\phi-\lambda\phi w_n}_{M_*}\longrightarrow0
  \quad\Longrightarrow\quad
  w_na-\beta^\phi_\lambda(a)w_n\longrightarrow0
  \quad *\text{-strongly}
\]
for every $\lambda>0$, every uniformly bounded sequence $(w_n)$ in $M$,
and every $a\in\BC(M,\phi)$.
\item\label{bc:ergodicity}
The fixed-point algebra of $\beta^\phi$ is $\C1$.
\end{enumerate}
\end{proposition}

\begin{proof}[Sources and hypotheses]
The restricted-state identity in \ref{bc:reduction} is
\cite[Proposition~3.4]{HI2017}; the type dichotomy and expectation are
\cite[Theorem~3.5]{HI2017}. The bicentralizer is globally invariant under
the modular group by \cite[Proposition~3.3]{HI2017}, so its canonical
expectation preserves $\phi$.
These results impose no amenability or fullness hypothesis.

For \ref{bc:flow}, apply \cite[Theorem~A(ii)]{AHHM2020} to the identity
inclusion $M\subset M$. Continuity is part of the flow construction.
The expectation-preservation identity in
\cite[Theorem~A(iv)]{AHHM2020} specializes to preservation of $\phi$,
because the relative commutant is $\C1$.

Finally, \cite[Theorem~D]{Marrakchi2020} states, for any $\sigma$-finite type
$\mathrm{III}_1$ factor, that every $\beta^\phi_\lambda$ with
$\lambda\ne1$ is ergodic. This is stronger than \ref{bc:ergodicity}.
Fullness is not assumed in this Theorem.
\end{proof}

\begin{proof}[Proof of \Cref{thm:main}, assuming \Cref{thm:rigidity}]
Suppose that $\BC(M,\phi)$ is nonscalar. Replace $(M,\phi)$ by
$(B,\phi|_B)$ using \Cref{bc:inputs}\ref{bc:reduction}.
The restricted state is faithful and normal.
The predual is still separable: if $E:M\to B$ is the normal expectation,
$B_*\ni\rho\mapsto\rho\circ E\in M_*$ is an isometric embedding.
We may therefore assume
\[
  \BC(M,\phi)=M,\qquad M\ne\C1.
\]

Every unitary $u\in M_\phi$ centralizes $\phi$ exactly. The unitary
definition of the bicentralizer gives $u^*au=a$ for every $a\in M$.
Since unitaries span $M_\phi$, and $M$ is a factor, it follows that
$M_\phi=\C1$.

Set $b_s=\beta^\phi_{e^{-s}}$.
The premise of \eqref{eq:flow-property} is equivalent to
\[
  \norm{w_n\phi-e^{-s}\phi w_n}_{M_*}\longrightarrow0.
\]
\Cref{bc:inputs}\ref{bc:flow} therefore gives exactly
\eqref{eq:flow-property}, and also continuity and state preservation.
By \Cref{thm:rigidity}, $b$ is trivial. By
\Cref{bc:inputs}\ref{bc:ergodicity}, it is ergodic. Hence $M=\C1$,
contradicting the reduction.
The contradiction applies to every original pair $(M,\phi)$.
\end{proof}
